\section{Introduction}

\subsection{Historical Background}
India’s agricultural heritage dates back to the Indus Valley Civilization, where sophisticated irrigation systems and granaries were already in use. Over centuries, agriculture in India evolved under various empires and rulers. During the colonial period, Indian agriculture suffered due to exploitative policies, leading to famines and rural poverty. However, after independence in 1947, significant efforts were made to revive and modernize the sector.

The most transformative event in India’s agricultural history was the Green Revolution in the 1960s and 70s by M. S. Swaminathan. With the introduction of high-yielding varieties (HYVs) of seeds, improved irrigation infrastructure, chemical fertilizers, and pesticides, food grain production, especially wheat and rice, saw a dramatic increase. It helped India achieve self-sufficiency in food production and reduced dependence on imports.

\subsection{Present-Day Scenario}
Today, agriculture in India is a vast and diverse sector. It includes the cultivation of crops, horticulture, dairy, fisheries, and forestry. The country is among the world’s top producers of rice, wheat, sugarcane, cotton, fruits, and vegetables. According to recent statistics, agriculture and allied sectors contribute around 15-18 \% to India’s Gross Domestic Product (GDP) and employ nearly 50 \% of the workforce.

India’s agriculture is heavily dependent on monsoons, making it vulnerable to climate variability. Over 50 \% of the net sown area is still rain-fed, despite an increase in irrigation coverage. In recent years, there has been a shift towards mechanization, organic farming, precision agriculture, and the use of agri-tech platforms to support farmers with weather forecasts, market prices, and modern techniques.

\subsection{About major crops in India }

Various food and non-food crops are grown in different parts of the country depending upon the variations in soil, climate, and cultivation practices.Major crops grown in India are rice, wheat, millet, pulses, tea, coffee, sugarcane, oil seeds, cotton, jute, etc. This can be classified into two classes:
\begin{itemize}
    \item \textbf{\textcolor{blue}{Food crops}}- Rice, Wheat, Millets, Maize and Pulses.
    \item \textbf{\textcolor{blue}{Cash crops}}- Sugarcane, Oilseeds, Horticulture crops, Tea, Coffee, Rubber, Cotton and Jute.
\end{itemize}
Also based on the season we can classify into three classes :- 

\begin{itemize}
    \item \textbf{\textcolor{blue}{Kharif crops}} [\textbf{Sown}: \textcolor{red}{June-July}
\textbf{Harvested}: \textcolor{red}{September-October} ]:- Rice (Aman), maize, jowar, bajra, tur, moong, urad, cotton, jute, groundnut, soybean etc.
\item \textbf{\textcolor{blue}{Rabi crops}} [\textbf{Sown:} \textcolor{red}{October-December}
\textbf{Harvested:} \textcolor{red}{April-June}]- Rice (boro) , Wheat, barley, peas, gram, mustard etc.
\item \textbf{\textcolor{blue}{Zaid crops}} [\textbf{Sown and harvested:} \textcolor{red}{March-July}
(between Rabi and Kharif)]- Rice(Aush) , Seasonal fruits, vegetables, fodder crops etc.
\end{itemize}
